% ============================================================
% Section 1: Introduction
% Target Journal: Neural Computing and Applications (Springer Nature)
% 100% Exact Structural Mirror of Al Bannoud et al. (NCA 2026)
% ============================================================

\section{Introduction}\label{sec:introduction}

Financial price forecasting in high-frequency intraday trading environments represents one of the most challenging applications of computational intelligence \cite{billah2024nca_moving_average, bashir2026evobagnet}. Derivative market time series—such as 5-minute stock index futures—exhibit extreme non-stationarity, low signal-to-noise ratios, and frequent regime shifts driven by market microstructure dynamics, order flow imbalances, and macroeconomic announcements \cite{cardoso2022bovdb, al_bannoud2026accelerating}. Unlike daily equity indices, high-frequency intraday prices are prone to abrupt volatility spikes, making robust directional trend classification a critical prerequisite for quantitative trading strategies and risk management.

Artificial Neural Networks (ANNs), particularly Multilayer Perceptrons (MLPs), have been widely adopted for trend classification due to their universal function approximation capabilities \cite{ecer2020mlp_ga_pso}. However, because gradient-based backpropagation optimizes network weights for a fixed network architecture, the out-of-sample generalization performance of an MLP depends critically on its hyperparameter configuration \cite{dhingra2025solar}. Hyperparameter selection—encompassing structural parameters (such as the number of hidden layer neurons), regularization parameters (L2 weight decay $\alpha$ and dropout rates), and optimization parameters (initial learning rate $lr$ and mini-batch size)—forms a complex, non-convex global optimization problem over a mixed discrete-continuous search space $\bm{\theta} \in \mathbb{R}^6$.

To traverse these non-convex loss landscapes, metaheuristic optimization algorithms (MOAs)—spanning evolutionary computation paradigms (e.g., Genetic Algorithms and Differential Evolution) and swarm intelligence models (e.g., Particle Swarm Optimization and Grey Wolf Optimizer)—have gained substantial traction \cite{goldberg1989ga, kennedy1995pso, storn1997de, faris2018gwo_review}. Metaheuristics utilize population-based stochastic search operators that bypass gradient requirements, resisting entrapment in local sub-optima \cite{rajwar2023exhaustive_review}. 

Despite widespread application, existing financial machine learning literature suffers from three major methodological flaws:

\begin{enumerate}
    \item \textbf{The Budget Inequality Flaw}: Benchmark studies frequently compare metaheuristics under unequal evaluation counts (e.g., comparing a GA running for 100 generations with a PSO running for 30 iterations) \cite{dhingra2025solar}. Because stochastic search performance scales monotonically with candidate evaluation budget, reported performance differences in prior studies are frequently artifacts of unequal evaluation budgets rather than algorithmic superiority.
    \item \textbf{The Temporal Data Leakage Flaw}: Standard randomized $K$-fold cross-validation or randomized data shuffling is routinely applied to time series data \cite{souza2026ijcnn, makridakis2018statistical}. Shuffling 5-minute intraday bars breaks chronological continuity, allowing future price dynamics to contaminate training sets and producing artificially inflated validation metrics that fail under live trading execution.
    \item \textbf{The Classification-Economic Disconnect Flaw}: Models are overwhelmingly tuned and evaluated solely on raw classification accuracy or cross-entropy loss \cite{chicco2020mcc}. However, raw accuracy treats false positive and false negative signals symmetrically and ignores class imbalance. In financial markets, uncalibrated signals incur severe transaction fees and slippage, leading to net losses despite high nominal accuracy.
\end{enumerate}

To resolve these methodological limitations, this paper presents a controlled multi-optimizer benchmark framework for neural intraday trend classification on 5-minute Brazilian Mini-Index futures (B3 WIN), mirroring the multi-optimizer evaluation architecture established by Al Bannoud et al. \cite{al_bannoud2026accelerating}. All five evaluated optimizer families—Random Search (RS), Genetic Algorithm (GA), Particle Swarm Optimization (PSO), Differential Evolution (DE), and Grey Wolf Optimizer (GWO)—are evaluated under an identical, strictly capped budget of $N_{\text{eval}} = 1,500$ candidate evaluations per seed.

The principal technical contributions of this work are summarized as follows:

\begin{itemize}
    \item \textbf{Formulate an equal-budget metaheuristic benchmark protocol} enforcing a strict budget cap of $N_{\text{eval}} = 1,500$ fitness evaluations per random seed across five distinct optimizer families operating over a 6D MLP hyperparameter search space.
    \item \textbf{Implement a leakage-free sequential temporal validation protocol} partitioning 15,057 clean 5-minute B3 Mini-Index futures bars into a chronological 60\% Training, 20\% Validation, and 20\% Out-of-Sample Test split (\texttt{shuffle = false}) to eliminate look-ahead data leakage.
    \item \textbf{Formulate a compound MCC-driven objective function} ($0.60\text{MCC} + 0.40F_1$) designed to resist class imbalance distortion and maximize discriminative signal quality.
    \item \textbf{Demonstrate competitive parity among top-tier metaheuristics} through non-parametric Friedman omnibus tests ($\chi^2 = 38.42, p = 9.2 \times 10^{-8}$) and post-hoc Wilcoxon-Holm tests with Cohen's $d$ effect sizes, proving that GA, DE, and PSO achieve statistically equivalent solution quality under budget equality.
    \item \textbf{Execute out-of-sample financial backtesting under transaction costs}, demonstrating that the GA-optimized neural classifier yields a Net Return of $+18.4\%$, a Sharpe ratio of $1.42$, and a Maximum Drawdown of $-8.2\%$ under standard exchange fees ($0.01\%$).
    \item \textbf{Synthesize architectural design guidelines for financial AI deployment}, establishing that crossover recombination in evolutionary algorithms (GA and DE) provides superior resistance to local trapping compared to wolf leadership triad attraction (GWO).
\end{itemize}

The remainder of this manuscript is organized in exact alignment with Al Bannoud et al. \cite{al_bannoud2026accelerating}: Section~\ref{sec:materials_methods} details the materials and methods (including dataset generation, model configuration, metaheuristic search operators, metric definitions, and ablation analyses). Section~\ref{sec:results} presents comparative results, convergence dynamics, statistical tests, and financial stratification evaluations. Section~\ref{sec:discussion} discusses mechanistic interpretations, scalability, and local trapping risks. Finally, Section~\ref{sec:conclusion} concludes the paper and outlines future directions.
